\chapter{Conclusions}
\label{ch:conclusion}

The population density increase in metropolitan areas, coupled with the higher pace of everyday life, calls for the development of Urban Air Mobility (UAM) solutions. 
To ensure sustainability and avoid exacerbating space constraints in cities, these vehicles must have a low ground footprint. 
The design of a Transwing eVTOL, which can offer a solution to this problem, constitutes the main objective of the project. 
This report aims to provide a final configuration for the vehicle that will be detailed further in the next stages of the design. This is done via a detailed trade-off, backed up by the preliminary design of the concepts.

After performing the trade-off, the Rotating Wing was chosen. This is a transformative design consisting of four engines attached to a wing that folds along the fuselage during take-off and landing. 
Its excellent performance in terms of "conversion effectiveness" made this a favourable option. 
This design would also perform relatively well on all the other aspects taken into account during the trade-off: TRL, transition efficiency, safety, and noise emissions. 
As far as the structural complexity is concerned, the bending moments that the hinge will need to withstand are predicted to be very high, however not hindering performance.

In order to evaluate the validity of the results, verification and validation procedures are carried out for all the preliminary sizing tools, which yielded positive results. 
A sensitivity analysis was performed on the trade-off to assess the influence of small errors in the evaluation and weighing of the criteria. 
The removal of an entire criteria and redistribution of the weights were considered. 
The Rotating Wing design proved to be the best of all options in 98\% of the possibilities, hence clearly showing the validity of the procedures. 

Despite the robustness demonstrated by the verification and validation procedures, and the sensitivity analysis, several limitations were encountered during the project. Firstly, the trade-off analysis relied heavily on preliminary data, which may not fully capture the complexities of real-world operations and environmental factors. 
The analysis was conducted using simplified models, which might not adequately reflect the true stresses and strains experienced during flight transitions, for example. 
Also, the assessment of transition complexity focused solely on the hinge, without considering the complexity of the actuators.
Additionally, the noise footprint analysis was based on theoretical projections, which need to be validated through empirical testing.
Finally, the assessment of the Technology Readiness Level (TRL) was somewhat speculative, as it did not account for potential delays or obstacles in the development process.

To build on the strong foundation established by the current project, several future adjustments and research directions are recommended.
First, a more comprehensive structural analysis should be performed, incorporating advanced simulation techniques and real-world testing to validate the preliminary findings. 
This will ensure the design can withstand the predicted stresses and enhance overall safety. 
Furthermore, noise footprint studies should transition from theoretical models to experimental setups, possibly using scaled prototypes to gather accurate data. 
Additionally, a detailed examination of the transition efficiency should be conducted under various operational conditions to optimise performance further. 
Moreover, an iterative design process should be adopted, allowing for continuous improvements based on testing outcomes. 
Finally, collaboration with regulatory bodies is essential to ensure the design meets all safety and operational standards, potentially streamlining the path from prototype to commercial implementation.

Following the successful selection of the Rotating Wing design, the next steps in the design process will focus on detailed engineering and optimisation. 
This phase will start by refining the aerodynamic properties of the design to enhance performance and stability during different phases of flight. 
Throughout this process, an iterative design approach will be employed, allowing for feedback and improvement. 
Ultimately, these efforts will culminate in the development of a small-scale prototype. 
The Rotating Wing Transwing eVTOL establishes the first steps to an innovative design with promising performance metrics. 
The eVTOL can provide a sustainable and space-efficient alternative for commuting within densely populated metropolitan areas and reducing traffic congestion and travel time. This is due to its ability to take off and land vertically on different types of landing sites, perfectly suitable for UAM.

%UAM, however, is only one of the possible applications of the vehicle, with industries like the one of Emergency Medical Services (EMS), logistics and parcel delivery, tourism and recreational services, and corporate and VIP transport being potentially revolutionised by this innovative solution. 

%The foundation of all successful design projects starts with the distribution of operational and technical roles and the establishment of project planning. Thus, after the construction of the project role hierarchy and definition of the project the project planning was performed. This was summarized in four different types of diagrams to allow easy understanding. These are the Project Organogram, the Work Flow Diagram (WFD) and the Work Breakdown Structure (WBS). Finally, the information from the WFD and WBS was transferred to a Gantt Chart.

%Subsequently, the baseline design was started by constructing the mission profile utilizing the requirements derived beforehand. Establishing the mission profile aided the procedure of the design option generation. The process of design discovery commenced with generating the design option tree (DOT) for the most essential subsystems. Three subsystems were deemed to bear the most importance based on the key requirements: the wing configuration, the ground footprint reduction and the propulsion configuration. After exploring the possibilities among these categories the Strawman Concept was used to eliminate the no-, infeasible and non-analysable concepts. The remaining options were shuffled together to create the final DOT on which the Strawman elimination was applied again.

%Before performing the trade-off of the designs a preliminary sizing of all four concepts was done. A mass estimation was performed based on historical data which served as an input to a more precise, semi-empirical, mass and power estimation model. After several design iterations, a final preliminary design was achieved for each concept. Furthermore, a verification plan of the tools used during the preliminary design was provided.

%Next, the trade-off was performed to achieve the final decision on the most suitable design concept. To accomplish this a sufficient number of trade-off criteria have been chosen. To cover the most important aspects the following six criteria were considered: Technology Readiness level, Transition efficiency,Energy Consumption, Structural complexity, Safety, and Noise Emission. Additionally, some criteria were further broken up into sub-criteria. A weight was assigned to each criterion and all concepts were assessed on a scale of 1 to 5. As a result the Rotating Wing (\ref{C2.1}) concept came out as winner. The result of the trade-off was reinforced by the sensitivity analysis.

%Finally, the design team assessed the interfaces in the design, the technical risk events, production and sustainability plan.

%This work aids the next steps of the design process and ensures the selection of the best possible design concept by the in-depth trade-off procedure. Furthermore, the established interfaces help to keep in mind the interdependencies between the subsystems during the entire design. The work shall be continued by further detailing the preliminary design of the Rotating Wing concept to obtain a more accurate preliminary mass and power budget. Afterwards, the detailed design of the vehicle can be commenced.
%\newpage
%\textit{This report provides an in-depth understanding of the baseline information %and key data and requirements based, on which the design of a VTOL vehicle used for %urban air mobility can be expanded and detailed.}
%
%As such, first and foremost, the purpose of this report was to showcase that the %design has a solid demand foundation, based on a performed market analysis. The %purpose of the market analysis was to identify the need for the product, leading to %the development of new requirements and changes in design where necessary. To %perform this task, a stakeholder analysis was performed, which was followed by a %marked segmentation analysis, a SWOT analysis, and a competition analysis.
%
%Based on the determined market segment, the mission profile of the vehicle was %determined and detailed. This mission profile was used in the development of a %functional analysis, which shall become necessary in the development and design of %the necessary systems and subsystems for the operation of the eVTOL. As such, a %functional flow diagram and a functional breakdown structure were provided to %strengthen the reader with a clear and easy-to-understand structure of the design %functions. Using this flow diagram, the hardware and software diagrams were then %conceived.
%
%Considering the functional analysis, as well as the needs of the stakeholder and the %project objective statement, a list of requirements and constraints has been %compiled and a requirement discovery tree has been developed, to enable the design %team to easily monitor the progress of the design, as well as compliance with all %the given requirements. Additionally, since sustainability has become a crucial %topic in recent years, the team has identified the importance of this aspect and as %such a sustainable development strategy has been conceived. Moreover, other crucial %aspects of the project such as risk and technical resources have been analysed in %detail, and contingency measures have been proposed to minimise detail.
%
%Lastly, this report has showcased a Design Option Tree that contains information %concerning all the configurations, which the team has considered conceivable. %Afterwards, based on feasibility and complexity criteria, some options have been %excluded, leaving five feasible design configurations for further analysis and trade-%off.
%
%Therefore, this report has provided a solid baseline for the future of the project, %coming up with a series of design configurations developed based on the mission %requirements imposed and the identified market segment.
